\subsection{A new ratio does not restart cumulative measurement}
\label{sec:repeated-item-rebalancing}

The origin of an added layer matters again at every later measurement.
Consider a continuing hedge of purchases with an unchanged hedging
instrument. The original purchase volume is 100 tonnes, designated at a
forward price of 80 currency units per tonne. For this illustration, assume
zero discounting and one common price risk; actual valuations must support
those simplifications. When the price reaches 90, the cumulative loss on
the hedged purchase is 1000. A separately supplied cumulative derivative
gain of 1200 gives reserve 1000 and cumulative ineffectiveness income 200.

Suppose the documented economic relationship now requires ten more tonnes
to be designated, while the risk-management objective remains unchanged.
Recognize the old designation's ineffectiveness first, then add those tonnes
at their current price of 90. At a later price of 95, the cumulative
hedged-item loss is
\[
 100(95-80)+10(95-90)=1550.
\]
If the cumulative derivative gain is then 1400, IFRS~9.6.5.11 requires a
reserve of $\min(1400,1550)=1400$, leaving cumulative ineffectiveness zero.
The derivative has gained only 200 since the last observation, but reserve
must increase by 400: debit derivative 200, debit ineffectiveness expense
200 and credit reserve 400. The earlier income of 200 reverses
\citep[IFRS~9, paragraphs 6.5.11(a)--(c), B6.5.8 and B6.5.17]{eu20231803}.

Restarting the lower-of calculation at the ratio change gives a different,
incorrect answer. It adds the smaller of the new derivative gain 200 and
new item loss 550 to the old reserve, leaving reserve $1000+200=1200$ and
profit 200. That journal can balance, yet it fails the cumulative rule.
Rebalancing continues the relationship; it does not turn successive
measurement intervals into separate hedges.

For repeated additions, retain each layer's quantity $q_\ell$ and
designation forward price $f_\ell$. At current price $p_t$, the signed
cumulative change in the hedged purchase is
\[
 H_t=-\sum_\ell q_\ell(p_t-f_\ell).
\]
The minus sign records a loss on the purchase when its price rises. Let
$D_t$ be the separately measured cumulative derivative gain. In this
finite offsetting case, the signed reserve and cumulative profit are
\[
 R_t=\operatorname{sgn}(D_t)\min(|D_t|,|H_t|),
 \qquad I_t=D_t-R_t.
\]
At each observation post only $\Delta D_t$, $-\Delta R_t$ and
$-\Delta I_t$ to derivative, reserve and gains, using debit-positive
signs. Their sum is zero because $\Delta D_t=\Delta R_t+\Delta I_t$.
Adding a layer at its current price gives that layer an initial change of
zero, so it creates no immediate journal and does not reset earlier amounts.
The same argument applies at each subsequent addition.

For example, add another 20 tonnes at 95 after the preceding measurement.
At price 97 the cumulative item loss is
$100(97-80)+10(97-90)+20(97-95)=1810$.
A cumulative derivative gain of 1900 then gives reserve 1810 and profit
90. Later price and derivative reversals are evaluated from these same
origins. They can reverse previously recognized ineffectiveness rather
than merely add another positive amount. Fewer intermediate observations
leave the final cumulative amounts unchanged when the designation history
and final valuations are the same.

The documentary \code{RepeatedItemRebalancing/1} report applies this
add-only item method to a finite dated event list. Each addition requires
an immediately preceding same-date measurement, a new layer identifier,
the current designation price and updated eligibility/documentation
findings. Missing findings are unresolved; different risk or valuation
methods and false eligibility findings are unsupported. Changes in units,
dates or layer origins reject the whole proposed report without posting.
The offset arithmetic does not itself establish an admissible hedge ratio.
Removed item volume, changed instrument volume and the reserve's treatment
when the purchase occurs require the separate rules described above.

